\chapter{Certified models and frontend comparisons}
\label{chap:certified-models}

The first four certificates use the legacy smooth-test domain. Two further certificates
identify ambient and intrinsic relaxations and transport the finite-dimensional legacy conclusions
to the selected closed-energy domain.  The tensorization normalization follows the classical
product argument in \cite{BakryGentilLedoux2014}.

\begin{definition}[Legacy smooth-test Poincar\'e API]
  \label{def:legacy-smooth-test-poincare}
  \lean{FunctionalInequalities.Poincare.Admissible,
        FunctionalInequalities.Poincare.Holds,
        FunctionalInequalities.Poincare.constant}
  \leanok
  For a Borel measure $\mu$ on a real Hilbert space, an admissible test is globally $C^\infty$,
  belongs to real $L^2(\mu)$, and has integrable squared canonical gradient.  The predicate
  $\operatorname{Holds}(\mu,C)$ asserts the corresponding Poincar\'e inequality, and
  $\CP(\mu)\in[0,\infty]$ is the infimum of valid finite nonnegative constants.
\end{definition}

\begin{theorem}[Linear-test and covariance lower bound]
  \label{cert:linear-test-lower}
  \lean{Solutions.LemLinearTestLower.certificate}
  \leanok
  \uses{def:legacy-smooth-test-poincare}
  If a probability law has finite second moment, every admissible Rayleigh quotient is at most
  $\CP(\mu)$, and $\lambda_{\max}(\Cov_\mu)\le\CP(\mu)$, including the zero-dimensional case.
\end{theorem}
\begin{proof}
  \leanok
  The certificate combines the supremum bound for admissible tests with the linear-test
  covariance computation and the Rayleigh characterization of its largest eigenvalue.
\end{proof}

\begin{theorem}[Binary tensorization]
  \label{cert:binary-tensorization}
  \lean{Solutions.LemPoincareTensorization.certificate}
  \leanok
  \uses{def:legacy-smooth-test-poincare}
  For probability laws $\mu$ and $\nu$ on separable real Hilbert spaces, fixed valid constants
  tensorize with $\max$, and
  \[
    \CP(\mu\otimes\nu)\le\max\{\CP(\mu),\CP(\nu)\}.
  \]
\end{theorem}
\begin{proof}
  \leanok
  Variance decomposition and almost-everywhere admissible sections give the fixed-constant
  inequality; taking the infimum yields the optimal upper bound.
\end{proof}

\begin{theorem}[Finite tensorization and equality]
  \label{cert:finite-tensorization}
  \lean{Solutions.LemPoincareFiniteTensorization.certificate}
  \leanok
  \uses{cert:binary-tensorization}
  For every finite, possibly dependently typed family of spaces and probability laws
  $(E_i,\mu_i)$, the product law---whose coordinates are independent by construction---satisfies
  \[
    \CP\!\left(\bigotimes_i\mu_i\right)=\sup_i\CP(\mu_i).
  \]
  The theorem includes empty, singleton, zero-dimensional, and infinite-constant boundary cases.
\end{theorem}
\begin{proof}
  \leanok
  Iterate the binary upper bound.  Coordinate lifts preserve variance and energy and give every
  reverse inequality.  The certificate packages the finite equality and all boundary laws.
\end{proof}

\begin{theorem}[Sharp one-dimensional Gaussian constant]
  \label{cert:gaussian-one-dimensional}
  \lean{Solutions.LemPoincareGaussianOneDimensional.certificate}
  \leanok
  \uses{cert:linear-test-lower,def:legacy-smooth-test-poincare}
  For every $m\in\R$ and every variance $v\ge0$, the Gaussian law
  $\gamma_{m,v}$ satisfies the legacy inequality at $v$, and
  \[
    \CP(\gamma_{m,v})=v.
  \]
  This includes the zero-variance Dirac law without a positivity premise.
\end{theorem}
\begin{proof}
  \leanok
  A direct one-dimensional density and Hardy estimate proves the upper bound for positive
  variance; the Dirac case is separate.  The identity test and the covariance lower bound give
  sharpness.
\end{proof}

The four model certificates use the legacy frontend. In finite dimension, the proved bridge
\cref{thm:legacy-canonical-constant-comparison} transfers their conclusions to the selected
constant without raw/core agreement. Tensorization and the Gaussian computation remain model
infrastructure: a general one-dimensional log-concave input is still required for all product
log-concave laws, and the long KLS routes require the remaining shared endpoint bridges.


\begin{theorem}[Certified ambient and intrinsic comparison]
  \label{cert:ambient-intrinsic-comparison}
  \lean{Solutions.LemAmbientIntrinsicPoincareComparison.certificate}
  \leanok
  \uses{lem:smooth-tests-affine-restriction-extension,
        thm:ambient-intrinsic-relaxation-equality,cor:ambient-intrinsic-constant-equality}
  The certificate packages smooth-test restriction and extension, equality of the two relaxed
  energies on the original $L^2$ carrier, and equality of their finite-constant predicates and
  extended optima for every finite-dimensional Borel probability.
\end{theorem}
\begin{proof}
  \leanok
  Combine the five geometric and optimizer comparisons without strengthening their hypotheses.
\end{proof}

\begin{theorem}[Certified legacy and relaxed comparison]
  \label{cert:legacy-relaxed-comparison}
  \lean{Solutions.LemLegacyRelaxedPoincareComparison.certificate}
  \leanok
  \uses{lem:preenergy-poincare-relaxation,thm:legacy-relaxed-holds-equivalence,
        thm:legacy-canonical-constant-comparison}
  The certificate packages the abstract pre-energy comparison for arbitrary measurable
  probabilities and the five finite-dimensional legacy-to-ambient-to-intrinsic comparisons.
  The Gaussian bound and exact optimum therefore transfer to the selected frontend, including
  zero variance. No locally Lipschitz or Cheeger identification is part of this certificate.
\end{theorem}
\begin{proof}
  \leanok
  Combine the two abstract comparisons with test-domain identification and the geometric bridge.
\end{proof}
